% Generated draft: requires mathematical and source review.
\documentclass[11pt]{article}
\usepackage[margin=25mm]{geometry}
\usepackage{fontspec}
\setmainfont{Latin Modern Roman}
\usepackage{amsmath,amssymb,mathtools,unicode-math}
\setmathfont{Latin Modern Math}
\usepackage{enumitem,needspace,xurl,hyperref}
\setlength{\parindent}{0pt}
\setlength{\parskip}{5pt}
\setlength{\emergencystretch}{3em}
\newcommand{\N}{\mathbb{N}}
\newcommand{\Z}{\mathbb{Z}}
\newcommand{\Q}{\mathbb{Q}}
\newcommand{\R}{\mathbb{R}}
\newcommand{\C}{\mathbb{C}}
\newcommand{\F}{\mathbb{F}}
\newcommand{\A}{\mathbb{A}}
\newcommand{\PP}{\mathbb{P}}
\newcommand{\E}{\mathbb{E}}
\newcommand{\eps}{\varepsilon}
\newcommand{\dd}{\,\mathrm{d}}
\begin{document}
\section*{UnsolvedMath ID: 7500007}
\textbf{Source problem code:} catalogue entry
\subsection{Definitions and mathematical statement}
View a logic $\mathcal L$ semantically: it assigns sentences to vocabularies and a satisfaction relation on structures. An expressive extension of first-order logic contains, or translates without changing truth, every first-order sentence. It is strictly stronger if some $\mathcal L$-sentence over some vocabulary has no first-order equivalent on all structures of that vocabulary.

Full compactness means that for every set $T$ of $\mathcal L$-sentences, if every finite subset of $T$ has a model, then $T$ has a model. Craig interpolation, in the sentence formulation, means: whenever $\varphi$ uses vocabulary $\tau_1$, $\psi$ uses $\tau_2$, and $\varphi\models\psi$, there exists an $\mathcal L$-sentence $\theta$ using only $\tau_1\cap\tau_2$ such that
\[
\varphi\models\theta\qquad\text{and}\qquad\theta\models\psi.
\]
The question asks whether a reasonable logic strictly extending first-order expressive power can satisfy both properties.

The adjective reasonable is not defined. Required closure under Boolean operations, quantification, substitution, vocabulary changes, isomorphisms or relativization, restrictions on syntax and languages, and any effectiveness conditions must be specified to make this an existence theorem in a fixed class of logics. The source also does not say whether compactness is full or only countable, or whether interpolation concerns sentences, formulas with parameters, or a stronger version. The definitions above make a standard full-compactness/sentence-interpolation candidate explicit, not an unnoticed choice of the missing reasonableness axioms. Strictly stronger concerns semantic expressive power, not merely a stronger proof calculus.
\subsection{Short English statement}
Is there a suitable abstract logic with more expressive power than first-order logic that retains compactness and Craig interpolation? The class of acceptable logics and precise versions of both properties need to be fixed.
\subsection{Sources}
\begin{itemize}
\item Local UnsolvedMath registry, permanent ID 7500007, exact source code catalogue entry, ``Compact Interpolation Logic Beyond First-Order Logic''; no proposer or year supplied. Registry status open, classification OPEN-TRIAGE.
\item MathOverflow, answer 227108, inventoried item 7, to ``What are some important but still unsolved problems in mathematical logic?'', \url{https://mathoverflow.net/questions/227083/what-are-some-important-but-still-unsolved-problems-in-mathematical-logic#227108}. Supplied direct-HTML secondary-source extraction, page/answer checked July 29, 2026. The post author's name is not included in the record, so complete author attribution cannot be supplied here.
\item The supplied summary says no reasonable stronger logic with both properties is known and no classification settles the question. It supplies no primary reference, axiom list for reasonable logics, or precise property variants; this is an unverified supplied assessment.
\item Provenance flags primary-source status NEEDS\_REVIEW. MathOverflow user contributions are CC BY-SA; post-author attribution and the applicable license version must be preserved, but the version is not given in the record.
\end{itemize}
\end{document}
